\chapter{Finite instruments and independently checkable access}
\label{finite:chapter}

\section{Representing an experiment without discarding its disturbance}

This chapter makes the abstract law profile constructive for finite
classical and quantum models. The algebraic representation is shared; the
conditions that make a representation physically valid differ. Throughout,
states are row vectors and effects are columns. A chronological product
$M_{s_1}M_{s_2}$ means that operation $s_1$ occurs first.

For candidate $h\in\{0,1\}$, a classical action $a$ has a finite recorded
outcome set $\mathcal Y_a$ and matrices $K_h^{a,y}$ on $n_h$ hidden states.
The entry $K_h^{a,y}(x,x')$ is the joint probability of recording $y$ and
ending at $x'$ conditional on starting at $x$. Physical validity requires
\begin{equation}\label{finite:classical}
 K_h^{a,y}\geq0\quad\hbox{entrywise},\qquad
 \sum_{y\in\mathcal Y_a}K_h^{a,y}\mathbf1=\mathbf1.
\end{equation}
The initial distribution $x_h$ is nonnegative and satisfies $x_h\mathbf1=1$.
An abort or failed preparation is a possible recorded outcome when it can
occur. Deleting its probability and renormalizing would define a different
experiment, usually with an unaccounted selection and preparation cost.

For a word of action/outcome pairs
$w=((a_1,y_1),\ldots,(a_m,y_m))$, the joint event probability is
\begin{equation}\label{finite:probability}
 p_h(w)=x_hK_h^{a_1,y_1}\cdots K_h^{a_m,y_m}\mathbf1.
\end{equation}
This expression retains both measurement probabilities and the change of
hidden state caused by measurement. An outcome-only noise model generally
cannot replace it when the detector disturbs subsequent dynamics.

For a quantum system of Hilbert dimension $q_h$, an action instead consists
of completely positive trace-nonincreasing maps
$\mathcal E_h^{a,y}$ whose sum is trace preserving. From density operator
$\rho_h$, the probability is
\begin{equation}
 p_h(w)=\operatorname{tr}\!\left[
 \mathcal E_h^{a_m,y_m}\circ\cdots\circ
 \mathcal E_h^{a_1,y_1}(\rho_h)\right].
\end{equation}
Expressing Hermitian operators in a real basis gives a row-state linear
representation with dimension $n_h=q_h^2$ and a terminal column $\tau_h$
representing trace. Transfer matrices can have negative entries in that
basis. Entrywise stochastic positivity is not a quantum validity test;
complete positivity and trace preservation must be established from the
instrument, for example through its Kraus operators.

Both cases therefore have the form
\begin{equation}\label{finite:linear}
 p_h(w)=x_hM_h^w\tau_h,\qquad
 M_h^w=M_h^{a_1,y_1}\cdots M_h^{a_m,y_m},
\end{equation}
on the common finite alphabet
$\Sigma=\{(a,y):y\in\mathcal Y_a\}$. A generic signed linear
representation is an algebraic input; it does not automatically define
probabilities. The exact software uses rational entries and recomputes
identities with rational arithmetic. The finite-dimensional theorems also
hold over the reals, but floating-point near-equality is not an exact
equivalence certificate.

\section{Finite closure decides equality of every finite record}

Pair two representations by setting
\begin{equation}\label{finite:paired}
 \alpha=(x_0,-x_1),\qquad
 \beta=\begin{pmatrix}\tau_0\\\tau_1\end{pmatrix},\qquad
 M_s=\begin{pmatrix}M_0^s&0\\0&M_1^s\end{pmatrix},\qquad
 d=n_0+n_1.
\end{equation}
Their contrast is $f(w)=\alpha M_w\beta=p_0(w)-p_1(w)$.
Let
\begin{equation}\label{finite:space}
 V=\operatorname{span}\{M_w\beta:w\in\Sigma^*\}.
\end{equation}
The generating columns are actual event effects; arbitrary linear
combinations in $V$ need not be effects of implementable measurements.
The span is a computational device for proving statements about the
generating experiments.

\begin{theorem}[Finite effect-space criterion]\label{finite:span}
For a finite alphabet and finite-dimensional representation, all contrasts
vanish if and only if $\alpha V=0$. If $\beta\ne0$ and some contrast is
nonzero, there is a distinguishing word of length at most $d-1$.
For valid instrument models this is equivalent to equality of the complete
laws of every fixed finite action sequence. Such equality also implies
equality under every common causal classical policy with almost surely
finite stopping under both models.
\end{theorem}
\begin{proof}
Let $V_k$ span effects of length at most $k$. Then
\begin{equation}
 V_0=\operatorname{span}\{\beta\},\qquad
 V_{k+1}=V_k+\sum_{s\in\Sigma}M_sV_k.
\end{equation}
If $V_{k+1}=V_k$, the space is invariant under all letters and contains
every word effect. Otherwise its dimension strictly increases. Starting
from dimension one, at most $d-1$ increases are possible. Hence
$V_{d-1}=V$, proving both the criterion and the witness bound. The case
$\beta=0$ has zero contrast for every word.

For valid instruments, words with a prescribed action sequence are the
atoms of its complete outcome law, so equality of all their probabilities
is precisely equality of that law. Under a causal randomized policy, a
complete action/outcome history has its instrument word probability
multiplied by the common probabilities with which the policy chooses its
actions and stopping along that history. These factors are identical under
both candidates. Equality therefore holds at each finite terminal history;
summing over the countable terminal histories proves equality for almost
surely finite stopping. Conversely, fixed action sequences are policies.
\end{proof}

Finite-span equivalence and short counterexample words are established
results for weighted and probabilistic automata
\cite{f_tzeng,f_kiefer,f_notes}. Observable operator models and predictive
state representations supply closely related descriptions of controlled
hidden dynamics \cite{f_thon,f_jaeger,f_littman}. Quantum equivalence also
has established finite-dimensional formulations \cite{f_wang}. The
framework uses these results as its exact algebraic foundation.

The allowed memory and control architecture must be represented. Arbitrary
unmodeled ancillas, coherent feedback, or additional continuous control
values are not included by calling the protocol ``adaptive.'' A specified
finite quantum memory can be included in the state dimension. Enlarging
that architecture changes the access problem.

\subsection{A certificate another implementation can check}

An equivalence certificate supplies columns $b_1,\ldots,b_r$, together with
coefficients establishing
\begin{equation}\label{finite:closure}
 \beta\in\operatorname{span}\{b_i\},\qquad
 M_s b_i\in\operatorname{span}\{b_j\}\quad\text{for all }s,i,
 \qquad \alpha b_i=0.
\end{equation}
Induction on word length places every $M_w\beta$ in the supplied span,
so every contrast vanishes. Linear independence of the columns is useful
for compact generation but unnecessary for sound checking. A non-equivalence
certificate supplies an actual word and its nonzero exact contrast. Neither
certificate requires the reviewer to trust the search execution.

For a minimal entry example, a one-state instrument emits a bit with
probability $p_h$ under action $a$. Its scalar kernels are $M_h^{a,1}=p_h$
and $M_h^{a,0}=1-p_h$. With $p_0=1/2$, $p_1=3/5$, the word $(a,1)$
has contrast $-1/10$, and the bound $d-1=1$ is attained. Both outcomes
remain possible under both hypotheses, so this exact witness is not a
perfect discriminator. Repeated sampling belongs to the statistical layer.

\section{Constraints on experiment histories}

Apparatus rules can prohibit particular action orders or require a
preparation and final readout. For an action-only deterministic finite
automaton, let $q_0$ be its initial state, $F$ its accepting states, and
$\delta(q,a)$ its possibly undefined transition. The matrix $A_a$ has entry
one at $(q,\delta(q,a))$ when that transition is defined and zero otherwise.
Let $v_F$ indicate accepting states. With $S$ automaton states, replace
\begin{equation}\label{finite:grammar}
 x_h\mapsto e_{q_0}^{\mathsf T}\otimes x_h,\qquad
 M_h^{a,y}\mapsto A_a\otimes M_h^{a,y},\qquad
 \tau_h\mapsto v_F\otimes\tau_h.
\end{equation}
The resulting series is the physical word probability on accepted complete
action sequences and zero elsewhere. It is a masked algebraic series,
not itself a normalized physical instrument on all strings. Applying
Theorem~\ref{finite:span} in paired dimension $Sd$ proves equality on
accepted experiments or produces an accepted witness of length at most
$Sd-1$. The lift is a standard finite-language construction, consistent
with constrained switching realization theory
\cite{f_petreczky,f_bastug}.

Nonaccepting intermediate grammar states are allowed. A protocol can pass
through one while continuing toward a legal completed experiment. Stopping
there is inadmissible. An empty accepted language gives a vacuous equality
certificate. For adaptive policies the supported terminal action sequences
must satisfy the grammar almost surely under both candidates.

Outcome-dependent restrictions require further care. A language that
retains only selected successful outcomes can describe postselection rather
than executable access. A physical specification must say what happens on
the excluded outcomes and how their probability and cost are retained.
The implemented grammar here is action-only; it does not silently solve
that more general policy-feasibility problem.

\section{The cheapest nonzero witness}

Give every action a nonnegative rational charge $c(a)$, independent of its
outcome and prior history, and define the lexicographic weight
\begin{equation}\label{finite:weight}
 \omega(w)=\left(c(w),|w|\right),\qquad
 c(w)=\sum_{j=1}^{|w|}c(a_j).
\end{equation}
The length component selects the shortest among equal-charge witnesses.
An action charge may be a calibrated upper bound on gross energy, a
specified duration, or a count. The theorem optimizes the supplied
additive quantity. A mean physical cost depending on the hidden state or
realized path generally does not have this form.

Use a priority queue ordered by $\omega$, with a deterministic final tie
breaker. Initially it contains $(\beta,\varnothing)$. Remove the least
candidate $(v,w)$. If $v$ is in the span of previously retained columns,
discard it. Otherwise retain the column and its actual word. Return $w$
if $\alpha v\ne0$. If the contrast is zero, insert every
$(M_s v,sw)$, where $sw$ prepends the letter $s$ to the chronological
word. This backward construction of effects does not reverse the
chronology of the resulting physical protocol.

\begin{theorem}[Cost-priority closure]\label{finite:cost}
The search accepts at most $d$ columns and removes at most
$1+|\Sigma|d$ candidates. It returns either an equivalence certificate or
a minimum-charge nonzero word, shortest among all minimum-charge nonzero
words. Zero-charge actions are allowed. The same result applies after
the grammar lift with its enlarged dimension.
\end{theorem}
\begin{proof}
Each retained column increases rank and only retained zero-contrast columns
generate candidates. This gives both finite bounds. In particular, the
search cannot loop through infinitely many zero-charge words.

A discarded candidate is a linear combination of earlier retained
columns. Since prepending increases length by one and adds a nonnegative
charge, queue removal weights never decrease. Every column in that
combination has weight no larger than the discarded candidate. Prepending
the same letter to the combination preserves the identity and the weight
ordering.

For completeness of this pruning rule, fix a weight threshold $T$ and
consider the state after all generated candidates of weight below $T$
have been processed, assuming none has nonzero contrast. Induct on word
length. The empty effect has been retained or expressed in retained
columns whenever its weight is below $T$. For $w=sv$ below $T$, the
inductive representation of $M_v\beta$ uses columns of weight at most
$\omega(v)$. Each of their $s$ prefixes was generated and has weight at
most $\omega(w)<T$. On processing, it was either retained or replaced by
earlier columns of no greater weight. Thus every below-threshold word
effect lies in the retained span, with no greater representative weight.

Before the returned witness, all retained effects have zero contrast. A
cheaper, or equal-charge shorter, nonzero word would have an effect in
their span, a contradiction. If the queue empties, the retained span
contains $\beta$ and every prefixed retained effect, and is annihilated
by $\alpha$. Equation~\eqref{finite:closure} then certifies equivalence.
\end{proof}

The returned word also gives a minimum-charge fixed action sequence whose
complete outcome laws differ: every outcome word for the same sequence has
the same action charge. It is an event to count while executing that action
sequence and retaining all outcomes. The witness does not confer a right
to postselect its outcomes without accounting for the failures.

This specialization is built on established finite-rank search. Its
historical priority is not established by this volume. It is useful because
it aligns a qualitative witness with a declared physical protocol charge.
It does not minimize repetitions, finite-error risk, total energy, or the
expected cost of adaptive stopping. A cheaper tiny gap can require much
greater total acquisition than a more expensive strong measurement.

\section{Certifying optimality without trusting the optimizer}

A reviewer should be able to verify more than the fact that a supplied
word distinguishes the candidates. The following certificate also proves
the declared minimum charge. Let $w_*$ be the proposed witness and
$T=\omega(w_*)$. Supply a finite list of actual-word effects
\begin{equation}
 b_i=M_{u_i}\beta,\qquad t_i=\omega(u_i)<T,
 \qquad\alpha b_i=0,
\end{equation}
and exact expansion coefficients with the following conditions:
\begin{align}
 \beta&=\sum_i r_i b_i,
 &r_i\ne0&\ \Longrightarrow\ t_i\leq(0,0),
 \label{finite:filteredbeta}\\
 M_s b_i&=\sum_j q_{sij}b_j,
 &q_{sij}\ne0&\ \Longrightarrow\
 t_j\leq t_i+(c(s),1).
 \label{finite:filteredclosure}
\end{align}
Require the initial expansion when $(0,0)<T$. Require the prefix expansion
for every $(s,i)$ with $t_i+(c(s),1)<T$. Recompute each effect, weight, and
contrast from the declared model and word. In particular, verify
$\alpha M_{w_*}\beta\ne0$.

\begin{theorem}[Filtered optimality certificate]
\label{finite:filtered}
A proposed nonzero word is minimum-charge and shortest among
minimum-charge nonzero words if and only if a finite certificate satisfying
these conditions exists. The certificate is sound with redundant columns,
signed matrices, and zero-charge actions.
\end{theorem}
\begin{proof}
For soundness, consider $w$ with $\omega(w)<T$. Induct on its length to
express $M_w\beta$ using only $b_i$ with $t_i\leq\omega(w)$. The empty
word uses Eq.~\eqref{finite:filteredbeta}. For $w=sv$, nonnegative
charges and the length tie-break give
$\omega(v)<\omega(sv)<T$. Every effect used for $v$ satisfies
\begin{equation}
 t_i+(c(s),1)\leq\omega(sv)<T.
\end{equation}
The corresponding required identity
\eqref{finite:filteredclosure} propagates the representation and its
weight bound. Since $\alpha b_i=0$, all below-threshold contrasts vanish.
The supplied nonzero witness therefore has the claimed optimum. When
$T=(0,0)$ there are no smaller words and no initial expansion is needed.

For completeness, run the cost-priority construction through all generated
candidates below $T$. If $w_*$ is optimal, every retained effect there
has zero contrast. Each processed discarded effect has an expansion in
earlier columns of no greater weight. The initial effect and every required
prefix were processed, giving the displayed identities. At most $d$
independent columns are needed, so this construction terminates even when
some actions have zero charge.
\end{proof}

For full equivalence, omit the threshold, require an expansion of $\beta$
and every prefix, and require zero contrast on all supplied columns.
This is both sound and complete by finite closure. The checker needs only
matrix multiplication, rational arithmetic, and the explicit expansion
identities; it need not calculate a basis or call an optimizer.

The weight restriction in the expansion is essential. Replacing a cheap
effect by an expensive representative without tracking its weight can
remove a required cheap prefix from the certificate and hide a cheaper
witness. Exact closure alone proves equivalence when imposed globally;
partial closure proves an optimum only when its filtration is respected.
The software tests deliberately forged word costs, missing prefixes,
model substitutions, and genuine but nonminimal witnesses.

A certificate must be bound to the model the reviewer intends to check.
A digest of an embedded model establishes internal consistency, but an
adversary could replace both model and digest. The checker therefore accepts
an independently selected expected model or digest. A digest is neither a
signature nor a calibration measurement. Physical validity and empirical
model adequacy remain separate requirements.

\section{What finite dimension does and does not make easy}

The dimension bound concerns the supplied representation. For two generic
$n$-qubit instruments, each real operator space has dimension $4^n$ and
the paired dimension is $2\cdot4^n$. Finite linear algebra therefore does
not imply a polynomial algorithm in qubit count, constant-shot process
characterization, or elimination of fault-tolerance overhead. Known
structure and a narrow target can reduce the represented problem, but
those are additional assumptions.

The number of rank operations is finite and polynomial in the explicit
matrix dimension for standard implementations. Exact arithmetic also incurs
bit complexity, and certificate size includes its rational coefficients.
A short word need not have a large probability contrast. Consequently a
bound on witness length is not a bound on the records required for reliable
inference. The statistical chapter keeps these quantities distinct.

\subsection{Parameterized families and an algebraic route}

Suppose every model has fixed representation dimension $n$, matrix entries
are rational functions of parameters $\theta$ on a semialgebraic domain,
and all denominators are nonzero there. Suppose the domain and functions
are effectively given over the rationals or algebraic numbers. For a
parameter pair $(\theta,\theta')$, Theorem~\ref{finite:span} makes
equivalence a finite collection of word-probability equalities for lengths
at most $2n-1$. Clearing nonvanishing denominators gives polynomial
conditions. For a rational target $g$, structural identifiability is the
validity of
\begin{equation}
 \forall\theta,\theta'\in\Theta:\quad
 \left(\bigwedge_{|w|\leq2n-1}p_\theta(w)=p_{\theta'}(w)\right)
 \Longrightarrow g(\theta)=g(\theta').
\end{equation}
Real quantifier elimination decides this statement in principle
\cite{f_basu,f_nemcova}. The finite conjunction can be enormous and the
decision can be computationally expensive. The observation is a route for
appropriately structured symbolic problems, not an implemented solver for
arbitrary continuous families. It also explains where algebraic geometry
enters naturally: the fibers of the observable map are polynomial or
rational constraint sets in such a model.

\section{An exact boundary for unbounded positive-gap search}

Deciding whether any word has a nonzero contrast is easier than deciding
whether some arbitrarily long experiment exceeds a fixed positive
discrimination threshold.

\begin{theorem}[Strict-cutpoint boundary in instrument form]
\label{finite:undecidable}
Fix a rational $\delta\in(0,1)$. There is no algorithm that, for every
pair of finite rational classical instruments with the following grammar,
decides whether some admissible finite protocol has transcript total
variation greater than $\delta$: any finite sequence of unobserved
stochastic actions, followed by exactly one terminal binary readout.
The horizon is unbounded and no repeated independent preparation is
available within a protocol.
\end{theorem}
\begin{proof}
Use the established undecidability of strict cutpoint emptiness for
rational probabilistic finite automata, in particular the fixed cutpoint
$1/7$ result described and proved in Ref.~\cite{f_rote}. Such an automaton
has initial distribution $\mu$, stochastic letter matrices $T_a$, and
binary acceptance column $v$. Its acceptance probability for word $w$ is
$p(w)=\mu T_wv$. The undecidable question is whether $p(w)>1/7$ for some
finite word.

Turn each letter into a single-outcome unobserved action. After these
actions allow one terminal readout with success effect
$D=aI+b\operatorname{diag}(v)$ and failure effect $I-D$. Choose rational
$b>0$ small enough that $a=\delta-b/7\geq0$ and $a+b\leq1$. These
diagonal readout kernels are physically valid classical instruments.
Compare against a one-state model that always reports zero. The only
informative observation is Bernoulli success with probability
$a+bp(w)$ versus certain failure; its total variation is $a+bp(w)$.
It exceeds $\delta$ exactly when $p(w)>1/7$.

All outcomes preceding the readout are deterministic, so feedback supplies
no information for choosing the word. If the chosen word is recorded,
randomization gives the average of its total variations, since the word
distribution is common to both models. If that average exceeds the
threshold, some deterministic word exceeds it. Thus an algorithm for
the stated instrument problem would decide strict cutpoint emptiness.
\end{proof}

This is a translation of an established computability result, historically
traced to Nasu and Honda, rather than a new impossibility theorem for
physics. The reduction does not show that its stochastic matrices are
embeddable continuous-time propagators, thermodynamically reversible, or
available in any particular laboratory. Allowing unlimited independent
resets and repeated measurements changes the problem: a nonzero single-run
gap can then be amplified statistically.

A finite horizon makes the specified finite-alphabet design problem
finite. A finite additive budget does so only when every repeatable action
has a strictly positive minimum charge, or a separate horizon is imposed.
Zero declared entropy production does not imply zero laboratory energy,
and a budget on one accounting channel need not bound the number of
possible operations. Bounded design, calibrated statistics, and physical
accounting are therefore required layers of the constructive framework.
